\section{Arbitrary hosts with a dominant maximum}\label{sec:dominant}

The maximum decomposition also gives a length-unrestricted family with
no restriction on the number of descents.

\begin{lemma}[Two words at each internal rank]\label{lem:twointernal}
If a host $w$ has at least two parts and has a part greater than one,
then $[z^k]F_w\ge2$ for $2\le k\le\wt w-1$.
\end{lemma}
\begin{proof}
Put $l=\len(w)$ and $n=\wt w$. If $k\le l$, use $1^k$ and a word
consisting of one two and $k-2$ ones, embedded using a non-one part.
If $k>l$, choose positive integers $c_i\le w_i$ with sum $k$.
If there are two such vectors, they give two distinct words. Otherwise,
choose a slack coordinate $j$, so $c_j<w_j$; this exists because $k<n$.
Every other $c_i$ equals one, since transferring a unit from $c_i>1$
to $c_j$ would give a second vector. Delete one of these other ones and
increase $c_j$ by one. This gives a second word of weight $k$, now of
length $l-1$.
\end{proof}

\begin{lemma}[A comparison at distance two]\label{lem:codimtwo}
For an arbitrary host of weight $n\ge5$, write $F_w=\sum p_i z^i$.
Then $p_{n-2}\ge p_3\ge p_2$.
\end{lemma}
\begin{proof}
The second inequality is Lemma~\ref{lem:half}. If $p_3\le2$, the first
follows from Lemma~\ref{lem:twointernal}, or directly when the host is
a single part or consists only of ones. Assume $p_3\ge3$. Write $e$
for the number of non-one parts and $h$ for the number of parts at
least three. Distinct words obtained by subtracting two units without
deleting a part number exactly $\binom e2+h$. These have the original
length; deleting one or two parts gives disjoint additional counts.

If $e\ge3$, this count is at least three. When $p_3=4$, the word $(3)$
occurs, so $h\ge1$, and the count is at least four.

Suppose $e=2$ and $h\ge1$. If both parts are at least three, the
same-length count is three. If a one is present, decrementing a
non-one part and deleting a one supplies another word. If no one is
present, $p_3\le3$, which is enough. If just one non-one part is at
least three, the same-length count is two. With no ones, deleting the
part two supplies a third word and $p_3\le3$. With ones present,
decrementing either non-one part and deleting a one gives two further
words: one retains two non-one parts and the other only one.

If $e=2$ and $h=0$, then $p_3\le3$ and at least one one is present.
The same-length count is one. If there are at least two runs of ones,
delete a one from each of two runs, and separately delete a part two;
these add two words of different lengths. If there is one run between
the two twos, deleting either two gives two distinct further words.
If the run is at an endpoint, reverse the host if needed and write it
as $1^r22$. Deleting a two gives $1^r2$, while deleting a one and
decrementing the final two gives $1^{r-1}21$. Thus the count is again
at least three.

It remains to consider $e=1$, with non-one part $t$.
If $t\ge3$ and the ones are all on one side, $p_3\ge3$ requires at
least two ones. Subtracting two from $t$, decrementing $t$ and deleting
a one, and deleting two ones give three different lengths; here
$p_3\le3$. If ones occur on both sides, subtracting two from $t$ and
deleting one from each side supply two words of different lengths.
Decrementing $t$ and deleting a one on either side supplies two more,
distinguished by the number of ones before the surviving non-one part.
Finally, if $t=2$ and $p_3\ge3$, ones occur on both sides. Since
$n\ge5$, one of the two runs has length at least two. Deleting two
ones from that run, or one from each run, gives two distinct words;
deleting the two gives a third word of a different length. Here
$p_3=3$. These cases exhaust the possibilities.
\end{proof}

\begin{lemma}[Central product inequality]\label{lem:centralproduct}
Let $A=F_X$, $B=F_Y$, $a=\wt X$, $b=\wt Y$, $D=a+b\ge6$, and
$C=AB=\sum c_i z^i$. Then
\[
 \Del A(D-1)+\Del B(D-1)+2c_{D-2}-c_{D-1}-c_3\ge0.
\]
Moreover $c_3\ge c_2$.
\end{lemma}
\begin{proof}
Interchange the arms so that $a\ge b$.
For a nonempty arm $Z$ of weight $n$, put
$q=p_{n-1}$, $r=p_{n-2}$ and $L=2r-p_2-p_3$.
Then $q\ge1$. Lemma~\ref{lem:codimtwo} gives $L\ge0$ for $n\ge5$.
For $n=4$ one has $L=r-q\ge1-q$, and for $n=3$ one has
$L=1-q$. Thus $L\ge1-q$ whenever $n\ge3$.

If $b=0$, the displayed expression is $r-p_3\ge0$ for the other arm.
If $b=1$, it is $L\ge0$, since $a\ge5$. If $b=2$, it is $L+q-1$;
here $a\ge4$, so this is nonnegative. If $b\ge3$, write $s,t$ for
the quantities $q,r$ of the second arm, and $L_B=2t-B_2-B_3$.
Neither arm contributes a boundary difference at rank $D-1$, and the
expression is
\[
 L+L_B+2qs-q-s
 \ge 2(q-1)(s-1)\ge0.
\]
These identities follow by expanding the top two and bottom three
coefficients of $AB$; all nonempty arm polynomials have constant and
linear coefficient one, and leading coefficient one.

For the last assertion, if one arm is empty use
Lemma~\ref{lem:half} for the other, whose weight is at least six.
Otherwise $c_3-c_2=A_3+B_3-1\ge0$, since at least one arm has
weight at least three.
\end{proof}

\begin{theorem}[Dominant maximum]\label{thm:dominant}
Let $w$ be an arbitrary nonempty composition of weight $N$ and maximum
part $M$. If $N\le2M+5$, then $F_w$ is unimodal.
\end{theorem}
\begin{proof}
The case $M=1$ is an interval polynomial. Assume $M\ge2$, split at
any maximum as $w=XMY$, and put $D=\wt X+\wt Y=N-M$.
The cases $D=0$ and a single empty arm cause no change to the identities
below. With $A=F_X$, $B=F_Y$, $C=AB$, the maximum decomposition gives
\begin{equation}\label{eq:dominantdifference}
 \Del F_w(n)=\Del A(n)+\Del B(n)
       +2c_{n-1}-c_n-c_{n-M-1}.
\end{equation}
Consequently all differences at $n\ge D+2$ are nonpositive.
If $M<D$, the difference at $D+1$ is also nonpositive:
when both arms are nonempty it equals $2-c_{D-M}\le0$, because
$c_j\ge2$ for $1\le j<D$; with an empty arm it equals
$1-c_{D-M}\le0$.

If $M\ge D-1$, Lemma~\ref{lem:half} gives increase through rank $D$,
so at most the single difference at $D+1$ remains undecided before
the permanently decreasing part. If $M=D-2$ or $M=D-3$, it gives
increase through rank $D-1$; the difference at $D$ is the only
undecided one. These cases are unimodal.

The two remaining cases are $M=D-4$ and $M=D-5$. Here $D\ge6$,
and first-half increase reaches $D-2$. At $n=D-1$, the subtracted
coefficient in~\eqref{eq:dominantdifference} is respectively $c_2$
or $c_3$. Lemma~\ref{lem:centralproduct} makes this difference
nonnegative in both cases. Again only the difference at $D$ remains
undecided, followed by nonpositive differences at every larger rank.
This proves unimodality.
\end{proof}

This theorem allows arbitrary lengths and descents. Its condition is
a sufficient one; no optimality of the constant five is asserted.
